% Originally: content/week-05.tex, \section{The Hessian test} (Corsin Week 5)

% Source: Corsin Nick/Class Notes/Week 5.pdf, p. 7
\section{The Hessian test}
\label{sec:hessian_test}

% Sonnet 5 (Medium): "Hessian" was already introduced in \cref{rem:second_order_taylor_hessian}
% (\cref{ch:taylor}); referencing it here instead of re-introducing the term with \newterm.
The \textbf{Hessian} (\cref{rem:second_order_taylor_hessian}) of $f \in C^2(U,\mathbb{R})$,
$U \subseteq \mathbb{R}^n$, is given by
\[\Hess f(x) = (\partial_i\partial_j f)_{i,j=1,\dots,n}\]
with respect to the standard basis of $\mathbb{R}^n$. By \newterm{Schwarz's lemma}
(\germanterm{Satz von Schwarz}), stated as \cref{thm:schwarz_mixed_partials},
$\Hess f(x)$ is symmetric.

% Supplement: Lukas Krause/Week 5 March 16-20 Notes Lukas Krause.pdf, pp. 2--3
\begin{ainote}
Symmetry of the Hessian is the hypothesis every statement in this section rests on --- the
spectral theorem in \cref{rem:diagonalizing_hessian}, the eigenvalue criteria, the Hurwitz
minors --- and neither Corsin nor \cref{thm:schwarz_mixed_partials} proves it. Lukas Krause's notes give a
complete proof, and unusually he leaves his \emph{failed} first attempt in place, with the
comment that it is instructive to see why the obvious calculation does not work. Both are
reproduced below, because the failure is genuinely the more informative half.
\end{ainote}

\begin{proof}[Proof of \cref{thm:schwarz_mixed_partials}]
It suffices to treat $m=1$ (apply the result to each component), $i<j$ (there is nothing to
prove when $i=j$), and $n=2$ with $i=1$, $j=2$ (freeze all the other variables). So let
$f \in C^2$ near $(x,y)$ and write $\partial_1,\partial_2$ for the two partial derivatives.

\textit{The attempt that fails.} Unwinding both mixed derivatives as iterated limits gives
\[\partial_2\partial_1f(x,y)
  = \lim_{h_2\to0}\lim_{h_1\to0}
    \frac{f(x+h_1,y+h_2) - f(x+h_1,y) - f(x,y+h_2) + f(x,y)}{h_1h_2},\]
and one would like to apply the mean value theorem in each variable and be done. The obstruction
is that the two intermediate points the mean value theorem returns depend on $h_1$ and $h_2$
\emph{separately} and need not agree, so forcing them together leaves a stray factor
$h_2/h_1$, which blows up when $h_1 \to 0$ at fixed $h_2$. The symmetric-looking expression
above is not enough on its own; one has to keep the two increments coupled.

\textit{The proof.} Couple them by fixing a single $h \neq 0$ and setting
\[F(h) := f(x+h,y+h) - f(x+h,y) - \big(f(x,y+h) - f(x,y)\big),\]
which is the second difference of $f$ over a square of side $h$, and
$\varphi(t) := f(x+th, y+h) - f(x+th, y)$, so that $F(h) = \varphi(1)-\varphi(0)$. The mean
value theorem gives $\xi_1 \in (0,1)$ with
\[F(h) = \varphi'(\xi_1) = h\,\partial_1f(x+\xi_1h,\ y+h) - h\,\partial_1f(x+\xi_1h,\ y).\]
Now apply the mean value theorem a second time, to $\psi(t) := \partial_1f(x+\xi_1h,\ y+th)$ ---
legitimate precisely because $f \in C^2$, so $\partial_1 f$ is itself continuously
differentiable in the second variable. Since $F(h) = h\big(\psi(1)-\psi(0)\big)$, there is
$\xi_2 \in (0,1)$ with
\[F(h) = h\,\psi'(\xi_2) = h^2\,\partial_2\partial_1f(x+\xi_1h,\ y+\xi_2h).\]
Grouping the four terms of $F(h)$ the other way --- first and third together, second and fourth
together --- runs the identical argument with the roles of the variables exchanged and produces
$\tilde\xi_1,\tilde\xi_2 \in (0,1)$ with
$F(h) = h^2\,\partial_1\partial_2f(x+\tilde\xi_1h,\ y+\tilde\xi_2h)$. Dividing both expressions
by $h^2 \neq 0$,
\[\partial_2\partial_1f(x+\xi_1h,\ y+\xi_2h)
  \ =\ \partial_1\partial_2f(x+\tilde\xi_1h,\ y+\tilde\xi_2h).\]
All four numbers $\xi_1,\xi_2,\tilde\xi_1,\tilde\xi_2$ lie in $(0,1)$ --- bounded, though they
depend on $h$ --- so both arguments converge to $(x,y)$ as $h \to 0$. Continuity of the mixed
second partials (again $f \in C^2$) lets us pass to the limit on both sides, giving
$\partial_2\partial_1f(x,y) = \partial_1\partial_2f(x,y)$.
\end{proof}

\begin{remark}[What the failed attempt shows]
\label{rem:why_schwarz_needs_care}
The two attempts differ in exactly one respect: the failed one lets $h_1$ and $h_2$ go to zero
\emph{independently}, the successful one ties them together as a single $h$ and only takes a
limit at the very end. That is why the second difference $F(h)$ over a \emph{square} is the
right object --- it is symmetric in the two variables by construction, so the two groupings of
its four terms compute the two mixed partials, and the symmetry of the answer is inherited from
the symmetry of the square rather than proved by force.

It also locates precisely where continuity is used: twice, once to apply the mean value theorem
to $\partial_1f$ and once to pass to the limit at the end. Drop it and the conclusion genuinely
fails --- see \cref{rem:schwarz_needs_c2} for the standard counterexample, which is
also the function Lukas sets as an exercise alongside this proof.
\end{remark}

\begin{definition}[Sign of a symmetric matrix]
\label{def:sign_symmetric_matrix}
A symmetric matrix $A \in \mathbb{R}^{n\times n}$ is:
\begin{itemize}
  \item \newterm{positive definite} (\germanterm{positiv definit}) if $v\transp Av > 0$ for all
    $v \in \mathbb{R}^n\setminus\{0\}$;
  \item \newterm{negative definite} (\germanterm{negativ definit}) if $v\transp Av < 0$ for all
    $v \in \mathbb{R}^n\setminus\{0\}$;
  \item \textbf{indefinite} if there exist $v, u \in \mathbb{R}^n\setminus\{0\}$ such that
    $v\transp Av > 0$ and $u\transp Au < 0$;
  \item \textbf{degenerate} (\germanterm{entartet}) if there exists
    $v \in \mathbb{R}^n\setminus\{0\}$ such that $Av = 0$; equivalently, if $\det A = 0$, or
    if $0$ is an eigenvalue of $A$.
\end{itemize}
\end{definition}

\begin{ainote}
Corsin defines degenerate as \qt{there exists $v \neq 0$ with $v\transp Av = 0$}. That
condition is much weaker than intended: if $A$ is indefinite, with $v\transp Av > 0$ and
$u\transp Au < 0$, then $t \mapsto (tv+(1-t)u)\transp A(tv+(1-t)u)$ is continuous and changes
sign, so it vanishes at some $w \neq 0$ --- making \emph{every} indefinite matrix degenerate,
and the third bullet of \cref{thm:hessian_test} vacuous. The intended notion is a nontrivial
kernel ($\det A = 0$), which is also what \cref{ex:6.2} states. Corrected above.

\end{ainote}

% Supplement: Sascha Brack/Class Notes/Week_05_Notes_Friday.pdf, pp. 3--4
\begin{remark}[Linear algebra interlude: diagonalizing the Hessian]
\label{rem:diagonalizing_hessian}
Since the Hessian $\Hess f(x_0)$ is a symmetric matrix, it is diagonalizable over $\mathbb{R}$ by the spectral theorem. There exist $n$ real eigenvalues $\lambda_1, \dots, \lambda_n \in \mathbb{R}$ and a corresponding orthonormal basis of eigenvectors $v_1, \dots, v_n \in \mathbb{R}^n$ such that $\Hess f(x_0) v_i = \lambda_i v_i$ for all $i$.

Every vector $x \in \mathbb{R}^n$ can be written in this basis as $x = \sum_{i=1}^n a_i v_i$ with $a_i \in \mathbb{R}$. Applying the Hessian yields:
\[
  \Hess f(x_0) x = \Hess f(x_0) \sum_{i=1}^n a_i v_i = \sum_{i=1}^n a_i \Hess f(x_0) v_i = \sum_{i=1}^n \lambda_i a_i v_i.
\]
In this eigenbasis, the Hessian matrix becomes diagonal. Consequently, the quadratic form evaluates to $x\transp \Hess f(x_0) x = \sum_{i=1}^n \lambda_i a_i^2$, which immediately links the definiteness of the Hessian to the signs of its eigenvalues.
\end{remark}

\begin{theorem}[The Hessian test]
\label{thm:hessian_test}
If $f \in C^2(U,\mathbb{R})$ has a critical point $x_0 \in \operatorname{int}(U)$, then:
\begin{itemize}
  \item $\Hess f(x_0)$ \textbf{positive definite} $\Rightarrow$ local \textbf{min} at $x_0$;
  \item $\Hess f(x_0)$ \textbf{negative definite} $\Rightarrow$ local \textbf{max} at $x_0$;
  \item $\Hess f(x_0)$ \textbf{indefinite} $\Rightarrow$ \newterm{saddle point}
    (\germanterm{Sattelpunkt}) at $x_0$.
\end{itemize}
If $\Hess f(x_0)$ is degenerate but not indefinite --- that is, positive or negative
semidefinite with $0$ an eigenvalue --- the test gives \textbf{no information}, and all three
behaviours are possible.
\end{theorem}

\begin{ainote}
Corsin's third bullet reads \qt{indefinite and not degenerate}. With degeneracy corrected to
$\det A = 0$ the extra hypothesis is unnecessary: indefiniteness alone already forces a saddle,
since directions $v$ with $v\transp\Hess f(x_0)v > 0$ and $u$ with $u\transp\Hess f(x_0)u < 0$
make $t\mapsto f(x_0+tv)$ increase and $t\mapsto f(x_0+tu)$ decrease away from $x_0$. Dropped
above, and the genuinely inconclusive case --- semidefinite \emph{and} degenerate --- stated
explicitly instead, since that is the one students actually run into
(\cref{ex:6.8} is exactly this trap).
\end{ainote}

% Moved here from the Lagrange-multiplier section (review item D2): it uses the Hessian,
% "indefinite" and the eigenvalue criterion, none of which were available at its former
% position 260 lines earlier.
\begin{aiexample}[Classifying saddle points via the Hessian]
\label{ex:ai_saddle_via_hessian}
% Generator: Gemini 3.6 Flash (Medium)
Consider $f(x,y) := x^2-y^2$ on $\mathbb{R}^2$. The gradient
$\nabla f(x,y) = (2x,\,-2y)\transp$ vanishes at $(0,0)$, and the Hessian is constant:
\[
\Hess f(0,0) = \begin{pmatrix} 2 & 0 \\ 0 & -2 \end{pmatrix}.
\]
Its eigenvalues are $\lambda_1 = 2 > 0$ and $\lambda_2 = -2 < 0$, so $\Hess f(0,0)$ is
indefinite and $(0,0)$ is a saddle point by \cref{thm:hessian_test} --- the model case the word
\qt{saddle} is named after, since $f$ rises along the $x$-axis and falls along the $y$-axis.
\end{aiexample}

\begin{remark}[Mnemonic]
positive $= \smile \dots$ min; negative $= \frown \dots$ max; indefinite $=$ saddle.
\end{remark}

\begin{center}
% Generator: Claude Opus 5 (High) --- FIG-W05-01: the three conclusive cases of
% \cref{thm:hessian_test}, drawn as the model quadratic each one is.
% Left:  x^2 + y^2, Hess = diag(2,2), both eigenvalues > 0 -- a bowl.
% Mid:  -x^2 - y^2, Hess = diag(-2,-2), both < 0 -- a dome.
% Right: x^2 - y^2, Hess = diag(2,-2), opposite signs -- a saddle (\cref{ex:ai_saddle_via_hessian}).
% In each panel the two coloured curves are the profiles along the two eigendirections; the
% signs of the eigenvalues are exactly the up/down curvature of those two profiles.
\begin{tikzpicture}[>=Latex, scale=1.0]

  % ---------- positive definite: bowl ----------
  \begin{scope}
    \draw[gray!55] (-1.5,0) -- (1.5,0);
    \draw[green!50!black, very thick, domain=-1.25:1.25, samples=60, smooth]
      plot (\x, {0.62*\x*\x});
    \draw[orange!90!black, very thick, domain=-1.25:1.25, samples=60, smooth]
      plot (\x, {0.30*\x*\x});
    \fill (0,0) circle (1.5pt);
    \node[font=\scriptsize, anchor=south west] at (0.06,0.02) {$x_0$};
    % Caption baseline shared with the other two panels so the row reads straight.
    \node[font=\small, align=center, anchor=north] at (0,-1.15)
      {\textbf{positive definite}\\ $\lambda_1,\lambda_2>0$\\ local \textbf{min}};
  \end{scope}

  % ---------- negative definite: dome ----------
  \begin{scope}[shift={(4.2,0)}]
    \draw[gray!55] (-1.5,0) -- (1.5,0);
    \draw[green!50!black, very thick, domain=-1.25:1.25, samples=60, smooth]
      plot (\x, {-0.62*\x*\x});
    \draw[orange!90!black, very thick, domain=-1.25:1.25, samples=60, smooth]
      plot (\x, {-0.30*\x*\x});
    \fill (0,0) circle (1.5pt);
    \node[font=\scriptsize, anchor=south west] at (0.06,0.06) {$x_0$};
    \node[font=\small, align=center, anchor=north] at (0,-1.15)
      {\textbf{negative definite}\\ $\lambda_1,\lambda_2<0$\\ local \textbf{max}};
  \end{scope}

  % ---------- indefinite: saddle ----------
  \begin{scope}[shift={(8.4,0)}]
    \draw[gray!55] (-1.5,0) -- (1.5,0);
    \draw[green!50!black, very thick, domain=-1.25:1.25, samples=60, smooth]
      plot (\x, {0.62*\x*\x});
    \draw[orange!90!black, very thick, domain=-1.25:1.25, samples=60, smooth]
      plot (\x, {-0.62*\x*\x});
    \fill (0,0) circle (1.5pt);
    \node[font=\scriptsize, anchor=south west] at (0.08,0.04) {$x_0$};
    \node[green!50!black, font=\scriptsize, anchor=west] at (1.28,0.90) {up};
    \node[orange!90!black, font=\scriptsize, anchor=west] at (1.28,-0.90) {down};
    \node[font=\small, align=center, anchor=north] at (0,-1.15)
      {\textbf{indefinite}\\ opposite signs\\ \textbf{saddle}};
  \end{scope}
\end{tikzpicture}
\end{center}

% Generator: Claude Opus 5 (High)
The picture is the whole theorem. By \cref{rem:diagonalizing_hessian} the second-order Taylor
term is $\tfrac12\sum_i\lambda_ia_i^2$ in the eigenbasis, so along the $i$-th eigendirection $f$
behaves like the parabola $\tfrac12\lambda_it^2$: it curves \emph{up} when $\lambda_i>0$ and
\emph{down} when $\lambda_i<0$. All up gives a minimum, all down a maximum, and one of each
gives a saddle --- there is a direction of ascent and a direction of descent through the same
point, so it is neither. And when some $\lambda_i = 0$ that direction is \emph{flat} to second
order, the quadratic term says nothing about it, and the test must fall silent --- which is
exactly the degenerate case, illustrated by $x^4+y^4$ against $x^4-y^4$ above.

\begin{ainote}
The regularity in \cref{thm:hessian_test} is stated by Corsin as $C^3$; the test needs only
$C^2$ (which is what the Hessian's definition assumes). Relaxed to $C^2$ above.
\end{ainote}

\begin{aiexample}[Inconclusive Hessian Test at a Degenerate Critical Point]
% Generator: Gemini 3.6 Flash (Medium)
Consider $f(x,y) := x^4 + y^4$ and $g(x,y) := x^4 - y^4$ at the critical point $(0,0)$.
\begin{itemize}
  \item For both functions, all second-order partial derivatives at $(0,0)$ vanish, so $\Hess f(0,0) = \Hess g(0,0) = \left(\begin{smallmatrix} 0 & 0 \\ 0 & 0 \end{smallmatrix}\right)$.
  \item The Hessian test is \textbf{inconclusive} since $\det \Hess = 0$.
  \item However, direct inspection reveals that $f(x,y) \geq 0 = f(0,0)$ for all $(x,y)$, so $(0,0)$ is a strict local minimum for $f$, whereas $g(x,0) > 0$ and $g(0,y) < 0$ for non-zero $x,y$, so $(0,0)$ is a saddle point for $g$.
\end{itemize}
\end{aiexample}

% Source: Corsin Nick/Class Notes/Week 5.pdf, pp. 8--9
\subsection{How do I determine positive/negative definiteness?}

Let $A \in \mathbb{R}^{n\times n}$ be symmetric with eigenvalues $\lambda_1,\dots,\lambda_n$.
Note: $\det(A) = \lambda_1\cdots\lambda_n$ and $\Tr(A) = \lambda_1+\dots+\lambda_n$.

\textbf{For $n = 2$:}
\begin{itemize}
  \item $\det(A) = \lambda_1\lambda_2 < 0 \implies$ \textbf{indefinite}
  \item $\det(A) > 0$, $\Tr(A) > 0 \implies$ \textbf{positive definite}
  \item $\det(A) > 0$, $\Tr(A) < 0 \implies$ \textbf{negative definite}
\end{itemize}

\textbf{For $n \geq 3$:}
\begin{itemize}
  \item Solve the characteristic equation $\det(A - \lambda\id) = 0$ for $\lambda$ and check all
    eigenvalues.
  \item \newterm{Hurwitz criterion} (\germanterm{Hurwitz-Kriterium}): for
    \[A = \begin{pmatrix} a_{11} & a_{12} & \cdots & a_{1n} \\ a_{21} & \ddots & & \vdots \\ \vdots & & \ddots & \\ a_{n1} & \cdots & & a_{nn}\end{pmatrix}\]
    define the leading principal minors
    \[A_1 = \det(a_{11}), \quad A_2 = \det\begin{pmatrix} a_{11} & a_{12} \\ a_{21} & a_{22}\end{pmatrix}, \quad \dots, \quad A_n = \det(A).\]
    Then:
    \begin{itemize}
      \item $A_1 > 0$, $A_2 > 0$, $A_3 > 0$, \ldots\ if and only if the matrix is \textbf{positive definite}
      \item $A_1 < 0$, $A_2 > 0$, $A_3 < 0$, \ldots\ if and only if the matrix is \textbf{negative definite}
      \item \textbf{indefinite} if neither $A_1 \geq 0, A_2 \geq 0, \dots$ nor
        $A_1 \leq 0, A_2 \geq 0, A_3 \leq 0, \dots$
    \end{itemize}
\end{itemize}

\begin{ainote}
Corsin writes the characteristic equation as ``$A - \lambda\id = 0$'', omitting the determinant.
Corrected above.
\end{ainote}

% Source: Corsin Nick/Class Notes/Week 5.pdf, pp. 9--10
\begin{example}[Definiteness calculations]
\begin{center}
\begin{tabular}{lll}
\textbf{Matrix} & \textbf{Computation} & \textbf{Conclusion} \\
\hline
$\left(\begin{smallmatrix}2&2\\2&5\end{smallmatrix}\right)$ & $\det = 6 > 0$, $\Tr = 7 > 0$ & positive definite \\
$\left(\begin{smallmatrix}-4&2\\2&-4\end{smallmatrix}\right)$ & $\det = 12 > 0$, $\Tr = -8 < 0$ & negative definite \\
$\left(\begin{smallmatrix}4&2\\2&-4\end{smallmatrix}\right)$ & $\det = -20 < 0$ & indefinite \\
$\left(\begin{smallmatrix}-5&0&0\\0&-4&-2\\0&-2&-4\end{smallmatrix}\right)$ & $A_1 = -5 < 0$, $A_2 = 20 > 0$, $A_3 = -60 < 0$ & negative definite \\
$\left(\begin{smallmatrix}-1&0&-4\\0&2&-2\\-4&-2&22\end{smallmatrix}\right)$ & $A_1 = -1 < 0$, $A_2 = -2 < 0$, $A_3 = -72 < 0$ & indefinite \\
\end{tabular}
\end{center}
\end{example}

% Source: Corsin Nick/Class Notes/Week 5.pdf, pp. 10--11
\begin{exercise}[Critical values of $x^3 - y^3 + 3\alpha xy$]
\label{ex:cubic_saddle_family}
Find the critical values of
\[f_\alpha : \mathbb{R}^2 \to \mathbb{R}, \qquad (x,y) \mapsto x^3 - y^3 + 3\alpha xy\]
for $\alpha \in \mathbb{R}\setminus\{0\}$ and determine min/max/saddle points.
\end{exercise}

% Kept inline: solved directly in Corsin Nick/Class Notes/Week 5.pdf, pp. 10--11.
\begin{exercisesolution}[Solution to \cref{ex:cubic_saddle_family}]
\[
\begin{aligned}
\frac{\partial f_\alpha}{\partial x} &= 3x^2 + 3\alpha y \overset{!}{=} 0 &&\implies y = -\frac{x^2}{\alpha} \\
\frac{\partial f_\alpha}{\partial y} &= -3y^2 + 3\alpha x \overset{!}{=} 0 &&\implies -\frac{x^4}{\alpha^2} + \alpha x = 0
\end{aligned}
\]
\begin{enumerate}[label=\textbf{\arabic*.}]
  \item $x = y = 0$
  \item $x \neq 0 \implies x^3 = \alpha^3 \implies x = \alpha$, $y = -\alpha$
\end{enumerate}
So, we have critical points $(0,0)$ and $(\alpha, -\alpha)$.

Compute the Hessian:
\[\begin{pmatrix} \partial_x^2 f & \partial_x\partial_y f \\ \partial_x\partial_y f & \partial_y^2 f\end{pmatrix} = \begin{pmatrix} 6x & 3\alpha \\ 3\alpha & -6y \end{pmatrix} = \Hess f(x,y)\]
\[\Hess f(0,0) = \begin{pmatrix} 0 & 3\alpha \\ 3\alpha & 0\end{pmatrix}, \qquad \det = -9\alpha^2 < 0 \implies \text{indefinite}\]
\begin{align*}
\Hess f(\alpha,-\alpha) &= \begin{pmatrix} 6\alpha & 3\alpha \\ 3\alpha & 6\alpha\end{pmatrix}
  = 3\alpha\begin{pmatrix} 2 & 1 \\ 1 & 2\end{pmatrix}, \\
\det &= 27\alpha^2 > 0, \qquad \Tr = 12\alpha.
\end{align*}
For $\alpha > 0$ this is positive definite; for $\alpha < 0$ it is negative definite. So:
\begin{itemize}
  \item $(0,0) \longrightarrow$ \textbf{saddle point}
  \item $(\alpha,-\alpha) \longrightarrow$ \textbf{local minimum} for $\alpha > 0$; \textbf{local
    maximum} for $\alpha < 0$
\end{itemize}
\end{exercisesolution}

% Originally: content/week-05.tex, end of the chapter

\begin{aiexercise}[Global Extrema on a Compact Triangle]
\label{ex:ai_extrema_triangle}
% Generator: Gemini 3.6 Flash (Medium)
Find the global maximum and minimum values of $f(x,y) := xy$ on the compact triangular region $T := \{(x,y) \in \mathbb{R}^2 \mid x \ge 0, y \ge 0, x+y \le 2\}$.
\end{aiexercise}


% --- exercises relocated here from the dissolved sheets ---

% Originally: content/week-06.tex, \exercisesheet{6}, problem 6.2
% Source: exercises/Ex6_Analysis2_eng.pdf

\begin{exercise}[6.2 --- The signature of a $2\times 2$ matrix \textnormal{(important)}]
\label{ex:6.2}
Despite the definition, it is not necessary to compute the eigenvalues of a matrix to find its
signature. Prove that for a symmetric $2\times 2$ matrix $M$ we have the following simple rule to
determine the signature in terms of $\det M$ and $\Tr M$:
\begin{itemize}
  \item If $\det M > 0$, $\Tr M > 0$ then $M$ is positive definite,
  \item If $\det M > 0$, $\Tr M < 0$ then $M$ is negative definite,
  \item If $\det M < 0$ then $M$ is indefinite,
  \item If $\det M = 0$, then $M$ is degenerate.
\end{itemize}
\end{exercise}

% Originally: content/week-06.tex, \exercisesheet{6}, problem 6.8
% Source: exercises/Ex6_Analysis2_eng.pdf

\begin{exercise}[6.8 --- Multiple choice \textnormal{(important)}]
\label{ex:6.8}
The Hessian matrix of $f \in C^2(\mathbb{R}^n)$ is positive semidefinite at a critical point $x_0$
of $f$, i.e.,
\[\langle v, \Hess f(x_0) v\rangle \geq 0 \quad \text{for all } v \in \mathbb{R}^n.\]
Which of the following statements necessarily hold? (There may be more than one).
\begin{enumerate}[label=\textbf{(\alph*)}]
  \item $x_0$ is a strict local minimum of $f$.
  \item $x_0$ is a local minimum of $f$.
  \item $x_0$ is not a local maximum of $f$.
  \item None of the above statements.
\end{enumerate}
\end{exercise}

% Originally: content/week-06.tex, \exercisesheet{6}, problem 6.9
% Source: exercises/Ex6_Analysis2_eng.pdf

\begin{exercise}[6.9 --- Minimization \textnormal{(important)}]
\label{ex:6.9}
The function $f : \mathbb{R}^2 \to \mathbb{R}$ is given by $f(x,y) = 2x^2+y^2-x$. Determine the
extrema of $f$ on\dots
\begin{enumerate}[label=\textbf{(\alph*)}]
  \item \dots the unit circle $S^1 = \{(x,y) \in \mathbb{R}^2 \mid x^2+y^2 = 1\}$;
  \item \dots the closed unit disk $D = \{(x,y) \in \mathbb{R}^2 \mid x^2+y^2 \leq 1\}$.
\end{enumerate}
\end{exercise}

% Originally: content/week-06.tex, \exercisesheet{6}, problem 6.12
% Source: exercises/Ex6_Analysis2_eng.pdf

\begin{exercise}[6.12 --- Critical points \textnormal{(important)}]
\label{ex:6.12}
Let $f : \mathbb{R}^2 \to \mathbb{R}$ be the function $f(x,y) = (ax^2+by^2)e^{-x^2-y^2}$ with real
parameters $a, b \in \mathbb{R}$. Find all critical points and determine their nature with the
Hessian test, depending on $a, b$.
\end{exercise}
